\chapter{Introduction and Scenario}

\section{Scenario Choice}

This project follows \textbf{Scenario C: Omnichannel Commerce Core}. VerdeMart Retail
started by using nopCommerce as its online storefront, but the business now needs more than a
web shop. It wants web sales, warehouse execution, shipping, store operations, support, and
customer visibility to work together as part of one commerce experience.

In this scenario, nopCommerce becomes the commerce core of that wider ecosystem. A web order
should be reflected in operational systems such as warehouse and shipping, and changes that
happen outside nopCommerce, such as stock or fulfillment updates, should become visible again
in the customer and operator experience. This is important because those surrounding systems
may be slow, stale, unavailable, or inconsistent, but the commerce experience still needs to
remain useful and understandable.

\section{Business Context and Drivers}

VerdeMart Retail began as a single-channel online shop. Three business pressures now demand an
architectural response:

\begin{enumerate}
    \item \textbf{Operational growth.} Fulfillment has moved to a dedicated warehouse, and
    physical stores now handle pickup and return flows. nopCommerce has no first-class model
    for communicating with either system, so operators manage these flows manually or through
    fragile workarounds.

    \item \textbf{Stock inconsistency.} The web catalog shows stock numbers from nopCommerce's
    local model. Warehouse and store sales can reduce physical stock without immediately
    updating the online view. This creates overselling risk and damages customer trust.

    \item \textbf{Operational fragility.} Any slow or unavailable external system currently
    causes visible degradation during checkout or shipment processing. There is no durable
    retry path, no degraded-state visibility, and no recovery workflow.
\end{enumerate}

These pressures translate directly into the quality attribute scenarios defined in Chapter 4:
availability during warehouse or shipping failure, consistency under stale inventory,
performance under campaign load, recoverability after provider outages, and modifiability as
new operational channels are added.

\section{Surrounding Systems in Scope}

The assignment permits the use of simulators or stubs for surrounding systems. For this
project, the following systems from the Scenario C ecosystem are represented:

\begin{itemize}
    \item \textbf{Warehouse execution system} (OpenBoxes or equivalent WMS): receives
    fulfillment requests, returns picked, packed, delayed, and failed acknowledgements.
    \item \textbf{Shipping provider} (WireMock simulator): handles label creation and
    tracking updates; can be made unavailable to test the recoverability scenario.
    \item \textbf{Store operations / POS}: publishes in-store sale events and pickup state
    changes that must update the online inventory view.
    \item \textbf{ERP / back-office} (ERPNext or equivalent): receives order and financial
    summary events from the commerce core.
\end{itemize}

RabbitMQ is used as the message broker for asynchronous propagation, consistent with the
Scenario C surrounding system list.

\section{Core Architectural Problem}

Currently, nopCommerce mainly acts as the online shop: it owns the storefront, catalog
browsing, cart, checkout, customer accounts, and order creation inside one web-centred
platform. This works while the business is mostly operating through the online channel, but it
does not address an omnichannel environment where warehouse execution, shipping, store
operations, and support systems also influence the customer experience.

The core architectural problem is deciding how this online shop can evolve into a commerce
core without assuming that every surrounding operational system is always fast, available, and
consistent. The architecture must define what remains inside nopCommerce, what moves around
it, and how the system behaves when warehouse or shipping state is delayed, stale, unavailable,
or later recovered.
